\documentclass[../../../include/open-logic-section]{subfiles}

\begin{document}

\olfileid[id]{sth}{replacement}{intro}
\olsection{Pendahuluan}

Penggantian merupakan skema aksioma yang membedakan $\ZF$ dari~$\Z$. Kita
telah menggunakannya begitu saja sepanjang
\crefrange{sth:ordinals::chap}{sth:spine::chap}. Dalam bab ini, akhirnya kita
akan mempertimbangkan pertanyaan berikut: apakah Penggantian dapat
dibenarkan?

Untuk mempertajam pertanyaan itu, patut diperhatikan bahwa Penggantian
sebenarnya cukup \emph{kuat}. Dalam bab ini (dan sekali lagi dalam
\olref[sth][card-arithmetic][fix]{sec}), kita akan memahami seberapa kuat
Penggantian itu. Namun, hal ini akan menyiratkan bahwa pembenaran memang
diperlukan.

Kita akan membahas dua jenis pembenaran. Secara garis besar: pembenaran
\emph{ekstrinsik} merupakan upaya membenarkan suatu aksioma berdasarkan
buahnya; pembenaran \emph{intrinsik} merupakan upaya membenarkan suatu
aksioma dengan menyarankan bahwa aksioma itu diteguhkan oleh konsep-konsep
matematis yang bersangkutan. Sepanjang bab ini kita akan memahami dengan
lebih baik apa artinya, tetapi pembahasan ini baru merupakan puncak gunung
es. Untuk pembahasan lebih lanjut, lihat terutama \citeauthor{Maddy1988a}
(\citeyear{Maddy1988a} dan \citeyear{Maddy1988b}).

\end{document}
